\section{Actual whole tensor morphism model norms}

\begin{lemma}[Original tensor morphism norm multiplicativity]
\label{lem:actual-model-tensor-hom-norm}
\lean{EnvelopeProperty.Berkovich.RelativeSchemeValuedPoint.modelHomNorm_tensorHom}
\leanok
\uses{thm:actual-model-line-tensor-norm,def:scheme-line-actual-model-hom-norm,lem:scheme-line-actual-model-norm-one-section}
For two actual original morphisms of actual invertible modules on
the same original proper model, the genuine model morphism norm
of their actual native tensor morphism equals the product of
their genuine original model morphism norms.
\end{lemma}
\begin{proof}
Choose an actual original generic section of norm one in each
source line. Their true pure tensor has norm one by the proved
actual model tensor norm theorem. Naturality of the genuine native
pullback tensorator and of the original pure tensor section says
that the actual image of this tensor by the pulled whole original
tensor morphism is the actual tensor of the two actual images.
The tensor norm theorem computes the norm of this image as the
product of the two original image norms. The original morphism
scaling theorem computes each as the corresponding model morphism
norm, since its source section has norm one. Applying the same
scaling theorem once to the whole actual tensor morphism and its
actual norm-one pure tensor identifies this same image norm with
the tensor morphism norm. These are the required equal original
values. No source section, tensor compatibility, or norm scaling
certificate is assumed.
\end{proof}
